\documentclass[12pt]{article}
\usepackage{amsmath, amssymb}
\usepackage{geometry}
\usepackage{hyperref}
\usepackage{physics}
\usepackage{cite}

\geometry{margin=1in}

\title{Filter Functions for Dynamical Decoupling Sequences}
\author{osipra}
\date{}

\begin{document}

\maketitle

\begin{abstract}
We derive explicit filter functions $F(\omega, t)$ for the Ramsey, Hahn echo, CPMG-$n$, XY-8, and UDD-$n$ pulse sequences. Each derivation follows from direct evaluation of the squared Fourier transform of the modulation function $f(t')$, as established in document 02. We characterize each filter function's spectral sensitivity and discuss its implications for noise spectroscopy.
\end{abstract}

\section{Ramsey Sequence}

For the Ramsey sequence there are no $\pi$ pulses, so $f(t') = 1$ for all $t' \in [0, t]$. The filter function integral evaluates to:

\begin{equation}
    \int_0^t e^{i\omega t'}\, dt' = \frac{1}{i\omega}\left(e^{i\omega t} - 1\right)
\end{equation}

Taking the magnitude squared:

\begin{equation}
    F_\text{Ramsey}(\omega, t) = \frac{1}{\omega^2}\left(e^{i\omega t} - 1\right)\left(e^{-i\omega t} - 1\right) = \frac{2 - 2\cos(\omega t)}{\omega^2}
\end{equation}

Using the identity $1 - \cos\theta = 2\sin^2(\theta/2)$:

\begin{equation}
    \boxed{F_\text{Ramsey}(\omega, t) = \frac{4\sin^2(\omega t/2)}{\omega^2}}
\end{equation}

In the limit $\omega \to 0$, using $\sin(\omega t/2) \approx \omega t/2$:

\begin{equation}
    F_\text{Ramsey}(\omega \to 0, t) = t^2
\end{equation}

The Ramsey filter function peaks at $\omega = 0$ with maximum value $t^2$. It is a \textbf{low-pass filter} — maximally sensitive to slow, DC-like noise and quasistatic fluctuations. This is why Ramsey dephasing in superconducting qubits is dominated by $1/f$ noise \cite{Bylander2011}.

\section{Hahn Echo Sequence}

For the Hahn echo, a single $\pi$ pulse is applied at $t/2$, so:

\begin{equation}
    f(t') = \begin{cases} +1 & 0 \leq t' < t/2 \\ -1 & t/2 \leq t' \leq t \end{cases}
\end{equation}

The filter function integral splits into two parts:

\begin{equation}
    \int_0^t f(t')\, e^{i\omega t'}\, dt' = \int_0^{t/2} e^{i\omega t'}\, dt' - \int_{t/2}^{t} e^{i\omega t'}\, dt'
\end{equation}

Evaluating each integral:

\begin{align}
    \int_0^{t/2} e^{i\omega t'}\, dt' &= \frac{1}{i\omega}\left(e^{i\omega t/2} - 1\right) \\
    \int_{t/2}^{t} e^{i\omega t'}\, dt' &= \frac{1}{i\omega}\left(e^{i\omega t} - e^{i\omega t/2}\right)
\end{align}

Combining:

\begin{equation}
    \frac{1}{i\omega}\left(2e^{i\omega t/2} - 1 - e^{i\omega t}\right)
\end{equation}

Factoring out $e^{i\omega t/2}$ and using $2 - e^{-i\omega t/2} - e^{i\omega t/2} = 2 - 2\cos(\omega t/2) = 4\sin^2(\omega t/4)$:

\begin{equation}
    = \frac{e^{i\omega t/2}}{i\omega}\cdot 4\sin^2\left(\frac{\omega t}{4}\right)
\end{equation}

Taking the magnitude squared (note $|e^{i\omega t/2}|^2 = 1$):

\begin{equation}
    \boxed{F_\text{Hahn}(\omega, t) = \frac{16\sin^4(\omega t/4)}{\omega^2}}
\end{equation}

In the limit $\omega \to 0$:

\begin{equation}
    F_\text{Hahn}(\omega \to 0, t) \propto \omega^2 t^4 \to 0
\end{equation}

The Hahn echo filter function \textbf{vanishes at $\omega = 0$}, confirming that the $\pi$ pulse refocuses slow noise completely. The filter peaks at:

\begin{equation}
    \frac{\omega t}{4} = \frac{\pi}{2} \implies \omega_\text{peak} = \frac{2\pi}{t}
\end{equation}

The Hahn echo is a \textbf{bandpass filter} centered at $\omega = 2\pi/t$. By varying $t$, the filter peak sweeps across the noise spectrum \cite{Cywinski2008}.

\section{CPMG-\texorpdfstring{$n$}{n} Sequence}

% TO BE DERIVED
%
% CPMG-n consists of n equally spaced pi pulses at times t_j = (2j-1)t/2n for j = 1,...,n.
% The modulation function f(t') alternates sign n times.
% The filter function peaks at omega = n*pi/t.
%
% Derivation to be filled in after working through by hand.

\section{XY-8 Sequence}

% TO BE DERIVED
%
% XY-8 is a phase-cycled version of CPMG with 8 pulses applied along alternating X and Y axes.
% The phase cycling suppresses pulse errors that accumulate in CPMG.
% The modulation function is identical to CPMG-8 for dephasing noise (Z-axis noise).
% Additional robustness to pulse errors comes from the phase alternation.
%
% Derivation to be filled in after working through by hand.

\section{UDD-\texorpdfstring{$n$}{n} Sequence}

% TO BE DERIVED
%
% UDD (Uhrig Dynamical Decoupling) places n pi pulses at nonuniform times:
%   t_j = t * sin^2(j*pi / 2(n+1)),  j = 1, ..., n
% This nonuniform spacing is optimal for suppressing dephasing to order n
% with only n pulses, for smooth S(omega).
%
% For sharp-featured S(omega) (e.g. TLS peaks), CPMG may outperform UDD.
%
% Derivation to be filled in after working through by hand.

\section{Summary}

\begin{table}[h]
\centering
\begin{tabular}{llll}
\hline
\textbf{Sequence} & \textbf{Filter function} & \textbf{Peak frequency} & \textbf{Character} \\
\hline
Ramsey      & $4\sin^2(\omega t/2)/\omega^2$   & $\omega = 0$       & Low-pass \\
Hahn echo   & $16\sin^4(\omega t/4)/\omega^2$  & $\omega = 2\pi/t$  & Bandpass \\
CPMG-$n$    & TBD                              & $\omega = n\pi/t$  & Tunable bandpass \\
XY-8        & TBD                              & $\omega = 8\pi/t$  & Phase-robust bandpass \\
UDD-$n$     & TBD                              & Nonuniform         & Optimal for smooth $S(\omega)$ \\
\hline
\end{tabular}
\caption{Filter functions for dynamical decoupling sequences used in this project.}
\end{table}

\begin{thebibliography}{9}

\bibitem{Cywinski2008}
L. Cywi\'{n}ski, R. M. Lutchyn, C. P. Nave, and S. Das Sarma,
\textit{How to enhance dephasing time in superconducting qubits},
Phys. Rev. B \textbf{77}, 174509 (2008).
\href{https://arxiv.org/abs/0807.4926}{arXiv:0807.4926}

\bibitem{Alvarez2011}
G. A. \'{A}lvarez and D. Suter,
\textit{Measuring the Spectrum of Colored Noise by Dynamical Decoupling},
Phys. Rev. Lett. \textbf{107}, 230501 (2011).

\bibitem{Bylander2011}
J. Bylander et al.,
\textit{Noise spectroscopy through dynamical decoupling with a superconducting flux qubit},
Nature Physics \textbf{7}, 565 (2011).

\bibitem{Uhrig2007}
G. S. Uhrig,
\textit{Keeping a quantum bit alive by optimized $\pi$-pulse sequences},
Phys. Rev. Lett. \textbf{98}, 100504 (2007).

\end{thebibliography}

\end{document}